%%
%% Chapter 5: Conclusion
%% Independent file, included by ../main.tex
%%

\section{Conclusion}
\label{sec:conclusion}

We presented the GT-conditioned CVAE, a conditional variational
extension of Gaussian-splatting-based sparse-view CBCT reconstruction.
The model encodes the ground-truth volume into the variational posterior
$q(z\mid C,Y)$ during training through a lightweight 3D CNN, yet
discards both the GT encoder and the dense views at deployment,
sampling exclusively from a sparse-view neural-process prior
$p(z\mid C)$. On the LUNA16 benchmark with six projections, the model
reaches $29.50$\,dB PSNR / $0.868$\,SSIM (best-eval), within $0.04$\,dB
of the deterministic backbone, while providing per-point predictive
variance at no additional inference cost.

The central methodological insight of this paper is that two design
choices carry the method. First, the decoder is conditioned on the
sparse context features, so the latent is the \emph{sole} channel
through which dense information reaches the volume during training;
this keeps the latent informative without any train/deploy mismatch.
Second, ablations show that \texttt{max}-fused anchor tokens preserve
the latent's discriminating signal, worth roughly $2$\,dB over
\texttt{mean} initialization. With these choices, the GT-conditioned
posterior concentrates the latent on case-specific content during
training at no cost at deployment.

\textbf{Future work} includes (i) validating the GT-conditioned design
under additional context counts and denser view schedules,
(ii) extending the model to hierarchical local latents for fine-grained
uncertainty maps, (iii) calibrating the predictive variance against
real reconstruction errors for use in dose-guided adaptive
radiotherapy, and (iv) transferring the information-flow design to other
modality-paired inverse problems (e.g., sparse-view MRI, low-dose PET).

\section*{Acknowledgements}

The authors gratefully acknowledge the computing resources and the
LUNA16 challenge data \cite{setio2017luna16}. The authors declare no
conflict of interest.